% !TEX root = ../main.tex
% =============================================================================
% CHAPTER 3 - RELATED WORK   (target ~4-5 pages)
% Tense: simple past for what a specific study did ("LaBraM introduced..."), present
% for claims that still hold. Always make the RELATION to your work explicit.
% =============================================================================

\chapter{Related Work}\label{ch:related}

Building on the concepts introduced in Chapter~\ref{ch:background}, this chapter reviews prior work most closely related to the thesis, focusing on clinical EEG classification, self-supervised EEG representation learning, spatial modelling, microstates, and generalisation. The chapter concludes by identifying the research gaps addressed in this work.


\section{Conventional EEG Analysis and Clinical Classification}\label{sec:rel_conventional}

Comparisons between feature-based and end to end approaches show that learned representations do not necessarily outperform well-designed handcrafted EEG features in clinical classification \parencite{craikDeepLearningElectroencephalogram2019,gemeinMachineLearningBasedDiagnostics2020}. This makes conventional feature-based methods an important reference point when evaluating whether self-supervised representations provide additional value.

\textcite{gemeinMachineLearningBasedDiagnostics2020} compared a comprehensive feature-based framework with end to end neural networks for distinguishing pathological from non-pathological EEG. Despite their different approaches to feature extraction, the feature-based and neural models achieved similar classification performance. Analyses of their decision patterns further indicated that both approaches relied on similar low-frequency information, particularly delta and theta activity at temporal electrodes. These findings demonstrate the importance of including strong handcrafted baselines when evaluating whether learned representations provide information beyond established EEG descriptors.

Supervised neural architectures provide an additional comparison with self-supervised pretrained representations. EEGNet \parencite{lawhernEEGNetCompactConvolutional2018} uses a compact convolutional architecture with depthwise and separable convolutions designed for EEG feature extraction. Building on EEGNet, EEGNeX \parencite{chenReliableSignalsDecoding2024} introduces additional convolutional processing, an inverted-bottleneck structure, and dilated convolutions to enhance spatial feature extraction and capture temporal dependencies over a larger receptive field. EEG Conformer \parencite{songEEGConformerConvolutional2023a} combines convolutional feature extraction with self-attention to model local features and longer-range dependencies. ATCNet \parencite{altaheriPhysicsInformedAttentionTemporal2023} combines convolutional processing, multi-head self-attention, and temporal convolutional networks within a sliding window framework to model temporal structure with a relatively small number of parameters. Although these architectures were evaluated mainly on brain-computer interface (BCI) and other task-based EEG paradigms, they provide useful supervised comparators for assessing whether self-supervised pretraining offers an advantage over representations learned directly from labelled data.


Clinical EEG studies also highlight the importance of evaluation design. DICE-Net \parencite{miltiadousDICENetNovelConvolution2023} combined relative band power and spectral coherence features with a convolutional-transformer architecture for dementia classification and evaluated performance using leave-one-subject-out validation. In schizophrenia classification, \textcite{shoeibiAutomaticDiagnosisSchizophrenia2021} reported high cross-validation performance using CNN-LSTM models on a cohort of 14 patients and 14 healthy controls. However, results obtained from such small cohorts provide limited evidence of generalisation to independent participants or datasets. More broadly, these studies illustrate that high internal classification performance alone is insufficient to establish robustness in clinical EEG, evaluation must also account for participant independence and, where possible, transfer to independently collected data.


\section{Self-Supervised Learning for EEG}\label{sec:rel_ssl}

Prior self-supervised EEG studies have investigated predictive, contrastive, and reconstruction based objectives \parencite{wengSelfSupervisedLearningElectroencephalogram2024}. Early clinical EEG work demonstrated that the design of the pretraining task strongly affects the information learned. \textcite{banvilleUncoveringStructureClinical2020} compared relative positioning, temporal shuffling, and contrastive predictive coding and evaluated the pretrained encoders with frozen linear classifiers. The learned features were particularly useful when few labels were available and captured information related to age, sleep stage, and pathology. However, the choice of negative examples also changed what the representation encoded: sampling negatives across recordings encouraged sensitivity to differences between recordings. This finding is important because self-supervised objectives can learn stable participant or recording characteristics as well as reusable physiological structure.

BENDR \parencite{kostasBENDRUsingTransformers2021} used a convolutional encoder and transformer to learn contextual representations by predicting latent future features among distractors. It was pretrained on a large clinical EEG corpus and transferred across several EEG tasks and recording setups. The results provided evidence that large-scale self-supervised pretraining can produce reusable EEG features, although most downstream tasks were not resting-state psychiatric or neurodegenerative classification.

Masked reconstruction provides a different form of self-supervision. MAEEG \parencite{chienMAEEGMaskedAutoencoder2022} masked latent EEG features and reconstructed the signal with a transformer-based decoder. In sleep staging experiments, more difficult masking configurations and longer input contexts produced stronger downstream representations than easier reconstruction tasks. EEG2Rep \parencite{foumaniEEG2RepEnhancingSelfSupervised2024} instead predicted latent target representations and used structured masking intended to preserve coherent context. Its evaluation included linear probing, finetuning, and transfer across EEG tasks. These studies suggest that masking ratio, mask geometry, input duration, and prediction target all influence the representations learned from EEG rather than acting as purely technical hyperparameters.

\section{Pretrained Self-Supervised EEG Models}\label{sec:rel_foundation}

Recent work has scaled self-supervised pretraining to larger and more heterogeneous EEG collections. These models differ in tokenisation, pretraining objectives, treatment of electrode information, and downstream adaptation \parencite{kuruppuEEGFoundationModels2025}. Their common goal is to learn transferable representations rather than training a new encoder from scratch for each task.

BrainLM \parencite{caroFOUNDATIONMODELBRAIN2024a} provides a closely related example from functional magnetic resonance imaging (fMRI) and forms the architectural basis of this thesis. It represents regional brain activity as temporal patches and uses masked autoencoding to reconstruct missing signal content from visible patches. BrainLM was evaluated on clinical variable prediction and external-cohort generalisation. However, fMRI and EEG differ substantially in their spatial and temporal characteristics, BrainLM operates on regional fMRI time series, whereas EEG consists of rapidly sampled signals recorded at scalp electrodes. Section~\ref{sec:m_brainlmeeg} describes how BrainLM is adapted to EEG in this thesis.

Several self-supervised EEG models address signal heterogeneity in different ways. BIOT \parencite{yangBIOTBiosignalTransformer2023} tokenises channels independently and combines spectral features with channel and temporal embeddings. LaBraM \parencite{jiangLARGEBRAINMODEL2024} learns discrete neural codes for EEG patches and predicts the codes of masked patches. EEGPT \parencite{wangEEGPTPretrainedTransformer2024} combines masked waveform reconstruction with spatiotemporal representation alignment and learned channel mappings. EEGFormer \parencite{chenEEGFormerTowardsTransferable2024} also learns discrete EEG representations and evaluates transfer from a large clinical corpus to an independent neonatal seizure dataset.

Other approaches focus more explicitly on spatial and temporal structure. CBraMod \parencite{wangCBraModCrissCrossBrain2025a} uses criss-cross attention to model spatial and temporal dependencies separately, while REVE \parencite{el2026reve} incorporates three-dimensional electrode coordinates and temporal position together with spatiotemporal masking. These approaches illustrate different assumptions about how EEG should be represented during pretraining. However, existing evaluations remain heterogeneous, and current evidence does not establish a single representation or pretraining strategy that is consistently superior across datasets and downstream tasks \parencite{kuruppuEEGFoundationModels2025}. This motivates controlled comparisons in which representation choices are varied while the underlying architecture, pretraining data, and downstream evaluation protocol are kept as consistent as possible.


\section{Spatial and Geometric Inductive Biases}\label{sec:rel_spatial}

Because EEG channels correspond to physical sensor locations, several approaches explicitly model relationships between electrodes. Graph-based methods represent electrodes as nodes and define connections using spatial distance, functional connectivity, or both. EEG-GCNN \parencite{waghEEGGCNNAugmentingElectroencephalogram2020} combined the geometry of a 10-20 montage with spectral coherence for neurological disease classification. The study demonstrated the potential value of spatial structure, although recordings from clinical and healthy groups were acquired using different systems, making acquisition differences a possible confound.

Other methods attempt to handle montage variation without defining a fixed graph. \textcite{yiLearningTopologyAgnosticEEG2023} mapped different electrode configurations to a common representation, combining multidimensional positional encoding with a channel hierarchy. BIOT \parencite{yangBIOTBiosignalTransformer2023} used learned channel embeddings, EEGPT \parencite{wangEEGPTPretrainedTransformer2024} used channel mappings, and REVE \parencite{el2026reve} derived positional information directly from electrode coordinates. These methods encode different assumptions: channel embeddings represent sensor identity, topology mappings align electrodes to a common structure, and coordinate-based methods preserve continuous physical location.

Rotary positional embedding \parencite[RoPE;][]{suRoFormerEnhancedTransformer2024}, introduced in Section~\ref{sec:bg_posenc}, provides another mechanism for representing position in Transformer attention. RoPE was originally developed for sequential positions rather than multidimensional electrode coordinates, so its application to EEG sensor geometry requires adapting the positional representation. Existing coordinate-based EEG models nevertheless suggest that explicit sensor geometry can be useful across varying electrode configurations \parencite{el2026reve}. The literature therefore motivates comparison between learned channel identity and geometry-aware spatial representations without assuming that physical proximity alone reflects neural connectivity.


\section{EEG Microstates in Clinical Populations}\label{sec:rel_microstates}

Prior microstate research, building on the analysis described in Section~\ref{sec:bg_microstates}, has examined both the clinical associations of these states and the reliability of their temporal properties \parencite{khannaMicrostatesRestingstateEEG2015}.

Clinical findings vary across disorders. A meta-analysis of schizophrenia reported increased occurrence of class C together with reduced duration and coverage of class D \parencite{rieger15YearsMicrostateResearch2016}. In mood and anxiety disorders, reported effects were smaller and more heterogeneous, including changes in the occurrence of classes B and D \parencite{chivuEEGMicrostatesMoodAnxiety2023}. These results indicate that microstates capture disease-related group differences, but they do not support interpreting individual microstate classes as diagnostic labels.

The reliability literature also shows that microstate results depend on how the templates are constructed. Using common maps across recordings produced more reliable estimates of duration, occurrence, and coverage than deriving an independent set of maps for every recording \parencite{khannaReliabilityRestingStateMicrostate2014a}. A larger test-retest study similarly found stronger reliability for duration, occurrence, and coverage than for transition measures \parencite{kleinertReliabilityEEGMicrostate2023}. Cross-study comparison is further complicated because the same class letter can refer to different scalp topographies; quantitative comparison of the maps is therefore more reliable than assuming that labels such as A, B, C, and D are identical across studies \parencite{koenigEEGMetaMicrostates2023}.

These findings make microstates relevant to representation learning as a compact description of recurring whole scalp organisation. At the same time, their dependence on template estimation and their variable clinical associations mean that they should be treated as structured neurophysiological information rather than fixed biological categories.

\section{Evaluation under Participant and Dataset Shift}\label{sec:rel_evaluation}

Evaluation practices are particularly important in EEG because one participant usually contributes many correlated recordings or windows. If data from the same participant appear in both training and test sets, performance can reflect stable participant characteristics rather than generalisation to unseen individuals. This issue was demonstrated in a clinical EEG archive where preventing repeated recordings from the same patients from crossing split boundaries reduced performance \parencite{kimDeepLearningBasedEEG2023}. Participant level validation, including leave-one-subject-out evaluation, therefore provides a stricter test of clinical generalisation \parencite{miltiadousDICENetNovelConvolution2023}.

For pretrained encoders, linear probing and finetuning also answer different questions. Frozen evaluation measures how directly task relevant information is available in the pretrained representation, whereas finetuning allows the encoder itself to adapt to the labelled task. Prior self-supervised EEG studies have used both approaches, and reviews report that finetuning can substantially outperform frozen evaluation \parencite{banvilleUncoveringStructureClinical2020,kuruppuEEGFoundationModels2025}. Dataset shift introduces a further challenge because EEG can vary in montage, hardware, sampling rate, filtering, reference, and noise characteristics. Benchmarking and robustness studies show that such changes can reduce both representation quality and predictive performance \parencite{engemannReusableBenchmarkBrainAge2021,waghEvaluatingLatentSpaceRobustness2022}. Cross-dataset evaluation is therefore needed when the intended claim extends beyond one acquisition setting.

\section{Research Gap}\label{sec:rel_positioning}

Prior work shows that handcrafted EEG features and supervised neural networks remain competitive baselines, while self-supervised pretraining can learn useful EEG representations \parencite{gemeinMachineLearningBasedDiagnostics2020,banvilleUncoveringStructureClinical2020,kostasBENDRUsingTransformers2021}. Recent pretrained models further demonstrate different approaches to masking, channel identity, temporal context, and electrode geometry \parencite{jiangLARGEBRAINMODEL2024,wangEEGPTPretrainedTransformer2024,wangCBraModCrissCrossBrain2025a,el2026reve}. However, evaluation varies substantially across tasks and protocols, and it remains unclear whether self-supervised representations provide consistent benefits for resting-state clinical EEG, particularly under participant-independent and cross-dataset evaluation \parencite{kuruppuEEGFoundationModels2025}.

Two related gaps concern representation design. First, coordinate-based spatial representations have been explored, but it remains unclear when explicit electrode geometry provides an advantage over learned channel identity \parencite{el2026reve}. Second, microstates capture recurring whole scalp organisation and show reproducible group level associations, but they have mainly been analysed as derived EEG features rather than incorporated into self-supervised representation learning \parencite{khannaReliabilityRestingStateMicrostate2014a,koenigEEGMetaMicrostates2023}. These gaps motivate a controlled comparison of self-supervised representations, spatial encoding, and microstate-informed pretraining, while separately evaluating participant level and cross-dataset generalisation.
